\documentclass[10pt,a4paper]{amsart}
\usepackage[margin=19mm]{geometry}
\usepackage{amsmath,amssymb,mathtools}
\usepackage[scr=rsfs,cal=euler]{mathalfa}
\usepackage{xfrac}
\usepackage[unicode,hidelinks,bookmarks=false]{hyperref}
\newcommand{\ie}{i.e.\ }
\newcommand{\cf}{cf.\ }
\newcommand{\resp}{respectively\ }
\newcommand{\Subsection}{\subsection}
\providecommand{\nobreakdash}{\nobreak\hbox{-}}
\setlength{\emergencystretch}{2em}
\multlinegap0pt
\newcommand{\Hyp}{\mathbb{H}}
\usepackage{amssymb}
\usepackage{mathtools}
\usepackage[shortlabels]{enumitem}
\newtheorem{theorem}{Theorem}
\newcommand{\C}{\mathsf{C}}
\newcommand{\F}{\mathsf{F}}
\newcommand{\J}{\mathcal{J}}
\title{Constant mean curvature hypersurfaces in $\Hyp^2\times\Hyp^2$ with double horocyclic symmetry}
\author{Julio Cesar Mohnsam}
\date{Source: arXiv:2607.00108v1 (CC BY 4.0)}
\begin{document}
\maketitle

\noindent\emph{Source and licence.} Constant mean curvature hypersurfaces in $\Hyp^2\times\Hyp^2$ with double horocyclic symmetry. Version arXiv:2607.00108v1 (CC BY 4.0), distributed by arXiv under the Creative Commons Attribution 4.0 International licence, http://creativecommons.org/licenses/by/4.0/, as recorded in the arXiv licence metadata for this exact version and shown on the abstract page https://arxiv.org/abs/2607.00108v1. Full licence notice: \texttt{licenses/arxiv-2607.00108-CC-BY-4.0.txt}.\par\smallskip

\noindent\textit{Editorial scope.} Counted unit: the article's theorem thm:integrais-reconstrucao (source0.tex lines 318--346). The article's complete original proof of it is delivered with it (source0.tex lines 348--398, one \texttt{proof} environment of 51 lines). No mathematics is written, repaired or re-derived: the statement and the proof are the article's own lines, in the article's own notation, and no retained line is substituted or re-flowed.  Retained uncounted as the source-local context the proof needs: lines 221--225; lines 292--298.  Nothing was omitted from any retained span, so the package carries no omission record.  No retained line contains a citation, so the wrapper carries no bibliography and no bibliography entry.  No retained line contains a figure environment, an image-inclusion command or drawing code, so no image is delivered and the package declares no asset dependency.  Editorial formatting only, in addition: an AMS \texttt{amsart} preamble that restates the article's own declarations and macros used by the retained lines, namely the environments \texttt{theorem} on separate counters, together with the standard packages those lines need.  The retained environments print in this document as Theorem 1 in order of appearance, whereas the article prints the same items with the numbering of its own full document.  The complete native source of the article as distributed in the arXiv e-print for this exact version is bundled unchanged as \texttt{sources/204147/source0.tex}.
\begin{equation}\label{eq:theta-def}
\frac{y'}{y}=\cos\theta,
\qquad
\frac{w'}{w}=\sin\theta,
\end{equation}

\begin{equation}\label{eq:edo-cmc}
\boxed{\;
\theta'=\F(\theta):=\sin\theta-\cos\theta+\C,
\qquad
\C=3H.
\;}
\end{equation}

\begin{theorem}[First integrals and universal reconstruction]
\label{thm:integrais-reconstrucao}
In every interval where $\F(\theta)\neq0$, the quantities
\begin{align}\label{eq:J1J2}
\J_1&=\frac{w}{y}\,e^{\C t-\theta},&
\J_2&=\frac{yw}{|\F(\theta)|}=\frac{yw}{|\theta'|}
\end{align}
are constant along any CMC solution. Furthermore, the first integrals $\J_1$ and $\J_2$ are functionally independent in the state space.
Consequently, if $\theta$ is a solution 
of~\eqref{eq:edo-cmc} with $\F(\theta)\neq0$, then
\begin{equation}\label{eq:yw-univ}
\boxed{\;
\begin{aligned}
y(t)&=\sqrt{\frac AB}\,\sqrt{|\F(\theta(t))|}\,
\exp\!\left[-\frac12\bigl(\theta(t)-\C t\bigr)\right]\\[2pt]
w(t)&=\sqrt{AB}\,\sqrt{|\F(\theta(t))|}\,
\exp\!\left[\frac12\bigl(\theta(t)-\C t\bigr)\right],
\end{aligned}
\;}
\end{equation}
is a solution $(y,w)$ of~\eqref{eq:theta-def} where
$A=\J_2$ and $B=\J_1$.
%, that is,
%\begin{equation}\label{eq:prod-razao}
%yw=A\,|\F(\theta)|,
%\qquad
%\frac wy=B\,e^{\theta-\C t}.
%\end{equation}
\end{theorem}

\begin{proof}
Let $\J_1$ and $\J_2$ be defined by~\eqref{eq:J1J2}.
Taking the logarithmic derivative and using~\eqref{eq:theta-def},
\eqref{eq:edo-cmc} and $\F'(\theta)=\cos\theta+\sin\theta$, we obtain
\begin{align*}
(\log\J_1)'
=
\frac{w'}w-\frac{y'}y+\C-\theta'
=
\sin\theta-\cos\theta+\C-\theta'
=
\F(\theta)-\theta'
=
0,
\end{align*}
and
\begin{align*}
(\log\J_2)'
=
\frac{y'}y+\frac{w'}w-\frac{\F'(\theta)\theta'}{\F(\theta)}
=
\cos\theta+\sin\theta-\F'(\theta)
=
0.
\end{align*}
Thus, $\J_1$ and $\J_2$ are constant.

For functional independence, let $\rho=\log y$ and $\sigma=\log w$.
In the state space with coordinates $(\rho,\sigma,\theta,t)$,
\begin{align*}
d\log\J_1
=
-d\rho+d\sigma-d\theta+\C\,dt,
\qquad
d\log\J_2
=
d\rho+d\sigma-\frac{\F'(\theta)}{\F(\theta)}\,d\theta.
\end{align*}
The components in $(d\rho,d\sigma)$ are $(-1,1)$ and $(1,1)$, whose
determinant is $-2\neq0$. Therefore, the differentials are linearly
independent.

Finally, setting $A=\J_2$ and $B=\J_1$, we obtain
\begin{align*}
yw=A|\F(\theta)|,
\qquad
\frac wy=B e^{\theta-\C t}.
\end{align*}
Multiplying and dividing these two identities, the formulas
\eqref{eq:yw-univ} follow.
\end{proof}

\end{document}
